% =====================================================================
%  Schnitthöhe: Das Blatt ragt nur knapp über das Werkstück hinaus.
%  Seitenansicht, Bedienperson rechts, Vorschub nach links.
%  Selbst gezeichnet (TikZ), Lizenz: CC BY-SA 3.0
% =====================================================================
\begin{tikzpicture}[x=1mm, y=1mm, font=\scriptsize]

  % Seitenansicht mit Werkstück der Dicke 8: \saegeseite{Blattüberstand}
  % (Tischoberkante bei y=14, Blattmitte bei x=40)
  \newcommand{\saegeseite}[1]{%
    \pgfmathsetmacro{\mitte}{14+8+#1-17}%
    \begin{scope}
      \clip (0,14) rectangle (82,46);
      \saegeblatt{40}{\mitte}{17}{24}
    \end{scope}
    % Spaltkeil: folgt dem Blatt im Abstand von etwa 2 mm
    \pgfmathsetmacro{\wa}{180-asin((14-\mitte)/19)}%
    \draw[keil] ($(40,\mitte)+(\wa:19)$) arc[start angle=\wa, end angle=112, radius=19]
      -- ($(40,\mitte)+(112:21.5)$) arc[start angle=112, end angle=\wa, radius=21.5] -- cycle;
    \fill[gehaeusegruen!70] (4,4) rectangle (78,12);
    \draw[tisch] (4,12) rectangle (78,14);
    \draw[werkstueck, fill opacity=0.6] (8,14) rectangle (76,22);
    \draw[vorschub] (70,26) -- (56,26);}

  % ---------------- 1 richtig --------------------------------------
  \feld{0}{0}{82}{52}{1\quad Blatt knapp über dem Werkstück}
  \saegeseite{3}
  \draw[mass] (60,22) -- (60,25) node[right=-0.5mm, midway] {};
  \draw[hilfslinie] (44,25) -- (62,25);
  \node[anchor=west, align=left] at (62,33) {knapp eine\\Zahnhöhe};
  \draw[hilfslinie] (61.5,30.5) -- (60,24);
  \pic at (74,45) {ok};
  \node at (41,7.5) {\color{white}wenig Blatt frei, sauberer Schnitt};

  % ---------------- 2 falsch ---------------------------------------
  \begin{scope}[shift={(88,0)}]
    \feld{0}{0}{82}{52}{2\quad Blatt zu hoch}
    \saegeseite{15}
    \pic at (74,45) {nok};
    \node at (41,7.5) {\color{white}viel freies Blatt, höhere Verletzungsgefahr};
  \end{scope}
\end{tikzpicture}
